\section{Homological invariants}
\label{sec:homological-invariants}

In this section, we determine the Ext groups between simple
$\Lambda_P$-modules and derive a combinatorial formula for the global
dimension of $\Lambda_P$.

\subsection{Ext groups between simple modules}
\label{subsec:simple-simple-ext-groups}

We now determine the Ext groups between the simple
$\Lambda_P$-modules.  By
\Cref{thm:standard-simple-ext-formula,cor:simple-costandard-ext-formula}, the standard--simple and
simple--costandard Ext groups associated with an interval $T$ are
simultaneously nonzero precisely for the intervals appearing in the
following family.
  For $S,C\in\Int(P)$, define
\begin{equation}
\label{eq:simple-simple-intermediate-family}
\mathcal W(S,C)
:=
\bigl\{
T\in\Int(P)
\bigm|
(S,T)\in\Upsilon^+(P)
\text{ and }
(T,C)\in\Upsilon^-(P)
\bigr\}.
\end{equation}
For $T\in\mathcal W(S,C)$, put
\begin{equation}
\label{eq:omega}
\omega(S,T,C)
:=
\omega^+(S,T)+\omega^-(T,C).
\end{equation}
We first show that $\mathcal W(S,C)$ forms a Boolean lattice whenever it is nonempty.

\begin{proposition}
\label{prop:intermediate-boolean-structure}
Suppose that $\mathcal W(S,C)$ is nonempty.  With respect to the inclusion order, it has a minimum
and a maximum given by
\begin{equation}
\label{eq:intermediate-extremal-intervals}
T_{\min}
=
S\cup C,
\qquad
T_{\max}
=
\conv
\{x \mid T\in\mathcal W(S,C), \ x \in T\}.
\end{equation}
Moreover,
\begin{equation}
\label{eq:intermediate-minimum-boundaries}
\partial_{T_{\min}}^+S
=
\Max(C\setminus S),
\qquad
\partial_{T_{\min}}^-C
=
\Min(S\setminus C),
\end{equation}
and
\begin{equation}
\label{eq:intermediate-maximum-boundaries}
\partial_{T_{\max}}^+S
=
\bigcup_{T\in\mathcal W(S,C)}\partial_T^+S,
\qquad
\partial_{T_{\max}}^-C
=
\bigcup_{T\in\mathcal W(S,C)}\partial_T^-C.
\end{equation}
The assignment
\begin{equation}
\label{eq:intermediate-boolean-isomorphism}
\mathcal W(S,C)
\longrightarrow
\bigl[
\partial_{T_{\min}}^+S,
\partial_{T_{\max}}^+S
\bigr]
\times
\bigl[
\partial_{T_{\min}}^-C,
\partial_{T_{\max}}^-C
\bigr],
\qquad
T
\longmapsto
\bigl(
\partial_T^+S,
\partial_T^-C
\bigr),
\end{equation}
is an isomorphism of posets, where the right-hand side is equipped with the product order induced by inclusion.
In particular, $\mathcal W(S,C)$ is a Boolean
lattice with rank function
\begin{equation}
\label{eq:intermediate-boolean-rank}
\operatorname{rk}_{S,C}(T)
=
\left|
\partial_T^+S
\setminus
\partial_{T_{\min}}^+S
\right|
+
\left|
\partial_T^-C
\setminus
\partial_{T_{\min}}^-C
\right|,
\end{equation}
and 
\begin{equation}
\label{eq:omega-boolean-rank}
\omega(S,T,C)
=
|\Max(C\setminus S)|
+
|\Min(S\setminus C)|
+
\operatorname{rk}_{S,C}(T)
\quad 
(T\in\mathcal W(S,C)).
\end{equation}
\end{proposition}


\begin{proof}
We first construct the maximal element.  Let
$T_1,T_2\in\mathcal W(S,C)$ and let $R$ be the convex hull of
$T_1\cup T_2$ in $P$.
The union $T_1\cup T_2$ is connected since both intervals contain
$S$, so $R$ is an interval.

Since $S$ is a relative downset in both $T_1$ and $T_2$, convexity
shows that $S$ is a relative downset in $R$.  Convexity also shows
that taking the convex hull creates no new cover from $S$, and hence
\[
\partial_R^+S
=
\partial_{T_1}^+S\cup\partial_{T_2}^+S
=
\Max(R\setminus S).
\]
Hence $(S,R)\in\Upsilon^+(P)$.
Dually, $(R,C)\in\Upsilon^-(P)$ with
\[
\partial_R^-C
=
\partial_{T_1}^-C\cup\partial_{T_2}^-C
=
\Min(R\setminus C).
\]
Therefore $R\in\mathcal W(S,C)$.

Since $\mathcal W(S,C)$ is finite, iterating this construction gives
the interval $T_{\max}$ in
\eqref{eq:intermediate-extremal-intervals}.  It contains every element
of $\mathcal W(S,C)$ and satisfies
\eqref{eq:intermediate-maximum-boundaries}.

Next, we characterize the elements of $\mathcal W(S,C)$ in terms of
$T_{\max}$.  We first record that every
$T\in\mathcal W(S,C)$ satisfies
\begin{equation}
\label{eq:intermediate-mandatory-boundaries}
\Max(C\setminus S)
\subseteq
\partial_T^+S,
\qquad
\Min(S\setminus C)
\subseteq
\partial_T^-C.
\end{equation}
Indeed, if $a\in\Max(C\setminus S)$ and $a<x$ for some
$x\in T\setminus S$, then $x\in C$ because $C$ is a relative upset
in $T$.  This contradicts the maximality of $a$ in $C\setminus S$.
The second inclusion is dual.



For $X$ and $Y$ satisfying
$\Max(C\setminus S)
\subseteq X \subseteq
\partial_{T_{\max}}^+S$ and
$\Min(S\setminus C) \subseteq Y \subseteq
\partial_{T_{\max}}^-C$,
define
\begin{equation}
\label{eq:intermediate-coordinate-inverse}
U(X,Y)
:=
\bigl(S\cup X\bigr)^{\downarrow_{T_{\max}}}
\cap
\bigl(C\cup Y\bigr)^{\uparrow_{T_{\max}}}.
\end{equation}
We claim that $U(X,Y)\in\mathcal W(S,C)$ and
\begin{equation}
\label{eq:intermediate-coordinate-boundaries}
\partial_{U(X,Y)}^+S
=
X,
\qquad
\partial_{U(X,Y)}^-C
=
Y.
\end{equation}

First, $S,C,X,Y\subseteq U(X,Y)$.  Indeed, every element of
$S\setminus C$ lies above an element of
$\Min(S\setminus C)\subseteq Y$, and dually every element of
$C\setminus S$ lies below an element of
$\Max(C\setminus S)\subseteq X$.  Hence
$S\subseteq(C\cup Y)^{\uparrow_{T_{\max}}}$ and
$C\subseteq(S\cup X)^{\downarrow_{T_{\max}}}$.
Since $X\subseteq\partial_{T_{\max}}^+S$ and
$Y\subseteq\partial_{T_{\max}}^-C$, it follows that $X$ and $Y$
also lie in both factors of
\eqref{eq:intermediate-coordinate-inverse}.

The set $U(X,Y)$ is convex.  Moreover, $S$ is a relative downset and
$C$ is a relative upset in $U(X,Y)$, since the same holds in
$T_{\max}$.  If $u\in U(X,Y)\setminus S$, then
$u\leq x$ for some $x\in X$.  Each $x\in X$ lies in
$\partial_{T_{\max}}^+S$, so $U(X,Y)$ is connected.  Thus
$U(X,Y)$ is an interval.

Every element of $U(X,Y)\setminus S$ lies below an element of $X$,
while
$X\subseteq\Max(T_{\max}\setminus S)$.  Hence
$\Max(U(X,Y)\setminus S)=X$.  Since $U(X,Y)$ is convex in
$T_{\max}$, a cover from $S$ to an element of $U(X,Y)$ is also a
cover in $T_{\max}$.  It follows that
$\partial_{U(X,Y)}^+S=X$.  Dually,
$\partial_{U(X,Y)}^-C=Y$.  Therefore
$U(X,Y)\in\mathcal W(S,C)$.

For every $T\in\mathcal W(S,C)$, we claim that
\[
T
=
U\bigl(
\partial_T^+S,
\partial_T^-C
\bigr).
\]
Indeed, every element of $T\setminus S$ lies below an element of
$\Max(T\setminus S)=\partial_T^+S$, and every element of
$T\setminus C$ lies above an element of
$\Min(T\setminus C)=\partial_T^-C$.
This gives the inclusion from left to right.
Conversely, every element of the right-hand side lies between two
elements of $T$, and hence belongs to $T$ by convexity.

Together with \eqref{eq:intermediate-coordinate-boundaries}, this shows
that the boundary map is a bijection
\[
\mathcal W(S,C)
\longrightarrow
\bigl[
\Max(C\setminus S),
\partial_{T_{\max}}^+S
\bigr]
\times
\bigl[
\Min(S\setminus C),
\partial_{T_{\max}}^-C
\bigr],
\qquad
T
\longmapsto
\bigl(
\partial_T^+S,
\partial_T^-C
\bigr),
\]
with inverse \eqref{eq:intermediate-coordinate-inverse}.
The boundary map preserves inclusion, since a cover between elements
of an interval is also a cover in $P$, while
\eqref{eq:intermediate-coordinate-inverse} is monotone in $X$ and
$Y$.  Hence this is an isomorphism of posets.

Finally, take
$X=\Max(C\setminus S)$ and 
$Y=\Min(S\setminus C)$.
By the preceding isomorphism, $U(X,Y)$ is the minimum element $T_{\min}$ of
$\mathcal W(S,C)$.  We now show that
$U(X,Y)=S\cup C$.
The inclusion $S\cup C\subseteq U(X,Y)$ is already known.
Suppose that
$x\in U(X,Y)\setminus(S\cup C)$.
Since $x\notin C$ and
$x\in(C\cup Y)^{\uparrow_{T_{\max}}}$, there exists
$y\in Y\subseteq S$ with $y\leq x$.
Along a saturated chain from $y$ to $x$, the first element outside
$S$ lies in
$\partial_{U(X,Y)}^+S=X\subseteq C$.
Since $C$ is a relative upset in $U(X,Y)$, this forces $x\in C$, a
contradiction.
Thus $U(X,Y)=S\cup C$.
Moreover,
\eqref{eq:intermediate-minimum-boundaries} follows from
\eqref{eq:intermediate-coordinate-boundaries}.


The two factors in
\eqref{eq:intermediate-boolean-isomorphism} are Boolean lattices
under inclusion. Their product is therefore a Boolean lattice, with
rank given by the sum of the two cardinality differences.
This is precisely
\eqref{eq:intermediate-boolean-rank}.
Since
$\omega^+(S,T)=|\partial_T^+S|$ and
$\omega^-(T,C)=|\partial_T^-C|$ for $T\in\mathcal W(S,C)$,
\eqref{eq:omega-boolean-rank} follows from
\eqref{eq:intermediate-minimum-boundaries} and
\eqref{eq:intermediate-boolean-rank}.
\end{proof}

We now return to the calculation of Ext groups. 
For $S,C\in\Int(P)$, define the graded $\kk$-vector space
$\mathsf K^\bullet(S,C)$ by
\begin{equation}
\label{eq:simple-simple-combinatorial-complex}
\mathsf K^p(S,C)
=
\bigoplus_{\substack{
T\in\mathcal W(S,C)\\
\omega(S,T,C)=p}}
\det(\partial_T^+S)
\otimes_\kk
\det(\partial_T^-C).
\end{equation}
For a cover $T\lessdot R$ in $\mathcal W(S,C)$,
\Cref{prop:intermediate-boolean-structure} shows that exactly one of
the two coordinates in
\eqref{eq:intermediate-boolean-isomorphism} changes.  Thus exactly one
of the following holds:
\begin{equation}
\label{eq:ss-upper-cover-data}
\partial_R^+S
=
\partial_T^+S\sqcup\{a\},
\qquad
\partial_R^-C
=
\partial_T^-C
\end{equation}
for a unique $a\in\partial_R^+S\setminus\partial_T^+S$, or
\begin{equation}
\label{eq:ss-lower-cover-data}
\partial_R^+S
=
\partial_T^+S,
\qquad
\partial_R^-C
=
\partial_T^-C\sqcup\{b\}
\end{equation}
for a unique $b\in\partial_R^-C\setminus\partial_T^-C$.
Define the corresponding component of the differential by
\begin{equation}
\label{eq:simple-simple-combinatorial-differential}
d^{\mathsf K}_{R,T}(\alpha\otimes\beta)
=
\begin{cases}
(a\wedge\alpha)\otimes\beta,
&
\text{in \eqref{eq:ss-upper-cover-data}},
\\[2mm]
(-1)^{|\partial_T^+S|}
\alpha\otimes(b\wedge\beta),
&
\text{in \eqref{eq:ss-lower-cover-data}}.
\end{cases}
\end{equation}
All other components are zero. 
The two length-two composites around each square in the Hasse diagram of $\mathcal W(S,C)$ cancel, so these maps define a differential
$d^{\mathsf K}$ on $\mathsf K^\bullet(S,C)$.


The following proposition identifies this combinatorial complex with
the derived Hom complex.  Its proof is given in
\Cref{app:simple-simple-assembly}.

\begin{proposition}
\label{prop:simple-simple-combinatorial-model}
For all $S,C\in\Int(P)$, there is an isomorphism
\begin{equation}
\label{eq:simple-simple-derived-combinatorial-model}
\bigl(\mathsf K^\bullet(S,C),d^{\mathsf K}\bigr)
\cong
\RHom_{\Lambda_P}\bigl(\LL(S),\LL(C)\bigr)
\end{equation}
in the derived category of cochain complexes of $\kk$-vector spaces.
\end{proposition}

The cohomology of $\mathsf K^\bullet(S,C)$ can now be read from
\Cref{prop:intermediate-boolean-structure}.  Suppose that
$\mathcal W(S,C)$ is nonempty.  By
\eqref{eq:intermediate-boolean-isomorphism} and
\eqref{eq:omega-boolean-rank}, if $\mathcal W(S,C)$ has more than one
element, then, up to a fixed one-dimensional factor and a
cohomological degree shift, $\mathsf K^\bullet(S,C)$ is the augmented
simplicial cochain complex of a nonempty simplex.  Hence it is
acyclic.  This leaves the case in which $\mathcal W(S,C)$ is a
singleton.

We call a pair $(S,C)\in\Int(P)^2$ a \emph{saturated pair} in $P$ if
$\mathcal W(S,C)$ is a singleton.  By
\Cref{prop:intermediate-boolean-structure}, in this case
$\mathcal W(S,C)=\{S\cup C\}$.  For such a pair, we simply write
\begin{equation}
\label{eq:saturated-pair-degree}
\bar\omega(S,C)
:=
\omega(S,S\cup C,C)
=
|\Max(C\setminus S)|
+
|\Min(S\setminus C)|.
\end{equation}
We denote by $\Sat(P)$ the set of all saturated pairs in $P$. 

Using \Cref{prop:simple-simple-combinatorial-model}, we prove the following. 

\begin{theorem}
\label{thm:simple-simple-ext-formula}
Let $S,C\in\Int(P)$ and $p\geq0$.  Then
\begin{equation}
\label{eq:simple-simple-Ext}
\Ext^p_{\Lambda_P}\bigl(\LL(S),\LL(C)\bigr)
\cong
\begin{cases}
\det\bigl(\Max(C\setminus S)\bigr)
\otimes_\kk
\det\bigl(\Min(S\setminus C)\bigr),
&
\text{$(S,C)\in \Sat(P)$, $p=\bar\omega(S,C)$},
\\
0,
&
\text{otherwise}.
\end{cases}
\end{equation}
\end{theorem}

\begin{proof}
By \Cref{prop:simple-simple-combinatorial-model}, the Ext groups are
the cohomology groups of $\mathsf K^\bullet(S,C)$.  If
$\mathcal W(S,C)$ is empty, the complex in
\eqref{eq:simple-simple-combinatorial-complex} is zero.  
If $\mathcal W(S,C)$ has more than one element, then
$(\mathsf K^\bullet(S,C),d^{\mathsf K})$ is acyclic by the preceding discussion.

Suppose that $(S,C)$ is saturated.  Then
\Cref{prop:intermediate-boolean-structure} gives
$\mathcal W(S,C)=\{S\cup C\}$.  Hence
\eqref{eq:simple-simple-combinatorial-complex} consists of a single
term.  By \eqref{eq:intermediate-minimum-boundaries}, this term is
\begin{equation}
\label{eq:saturated-pair-determinant-line}
\det\bigl(\Max(C\setminus S)\bigr)
\otimes_\kk
\det\bigl(\Min(S\setminus C)\bigr),
\end{equation}
and by \eqref{eq:saturated-pair-degree} it occurs in degree
$\bar\omega(S,C)$.  This proves
\eqref{eq:simple-simple-Ext}.
\end{proof}



For $S,C\in\Int(P)$ and $p\geq0$, let
$\beta_{\LL(S)}^p(C)$ denote the $p$th Betti multiplicity of $\LL(S)$ at $C$, that is, the multiplicity of $\PP(C)$ in the
$p$th term of a minimal projective resolution of $\LL(S)$.

\begin{corollary}
\label{cor:simple-simple-Betti-projective-dimension}
For $S,C\in\Int(P)$ and $p\geq0$,
\begin{equation}
\label{eq:simple-simple-Betti-multiplicities}
\beta_{\LL(S)}^p(C)
=
\begin{cases}
1,
&
\text{$(S,C)\in \Sat(P)$ and $p=\bar\omega(S,C)$},
\\
0
&
\text{otherwise}.
\end{cases}
\end{equation}
In particular, every indecomposable projective occurs with multiplicity
at most one in the minimal projective resolution of $\LL(S)$, and both
the multiplicities and the degrees in which they occur are independent
of the coefficient field.
\end{corollary}


Consequently,
\begin{equation}
\label{eq:simple-simple-projective-dimension}
\pd_{\Lambda_P}\LL(S)
=
\max_{\substack{
C\in\Int(P)\\
(S,C)\in \Sat(P)
}}
\bar\omega(S,C).
\end{equation}


\subsection{Global dimension}
\label{sec:global-dimension}

The projective dimension formula in \eqref{eq:simple-simple-projective-dimension} gives an exact combinatorial expression
for the global dimension of $\Lambda_P$.  Define
\begin{equation}
\label{eq:global-dimension-combinatorial-invariant}
\Omega(P)
:=
\max_{(S,C)\in\Sat(P)}
\bar\omega(S,C).
\end{equation}

\begin{theorem}
\label{thm:global-dimension-formula}
For every finite connected poset $P$,
\begin{equation}
\label{eq:global-dimension-formula}
\gldim\Lambda_P=\Omega(P).
\end{equation}
In particular, the global dimension of $\Lambda_P$ is independent of the coefficient field $\kk$.
\end{theorem}

\begin{proof}
Taking the maximum of
\eqref{eq:simple-simple-projective-dimension} over
$S\in\Int(P)$ gives
\eqref{eq:global-dimension-formula}.
\end{proof}


The formula for $\Omega(P)$ also yields a sharp uniform upper bound.

\begin{corollary}    
\label{cor:uniform-global-dimension-bound}
For $n\geq3$, every finite connected poset $P$ with $|P|=n$
satisfies
\begin{equation}
\label{eq:uniform-global-dimension-bound}
\gldim\Lambda_P\leq n-1.
\end{equation}
For each $n \geq 3$, the bound in
\eqref{eq:uniform-global-dimension-bound} is attained by some such
poset.
\end{corollary}




\begin{proof}
Let $(S,C)\in\Sat(P)$.  Then
\[
\bar\omega(S,C)
=
|\Max(C\setminus S)|
+
|\Min(S\setminus C)|
\leq n.
\]
If equality held, then $S\cap C=\varnothing$,
$S\cup C=P$, $\Max(C)=C$, and $\Min(S)=S$.
Since $S$ and $C$ are connected, both would be singletons, which
contradicts $n\geq3$.  Hence
$\bar\omega(S,C)\leq n-1$ for every $(S,C)\in\Sat(P)$.
Thus $\Omega(P)\leq n-1$, and
\Cref{thm:global-dimension-formula} gives
\eqref{eq:uniform-global-dimension-bound}.

For sharpness, fix $n\geq3$ and consider the poset
$P_n=\{s,u_1,\ldots,u_{n-1}\}$,
where $s<u_i$ for $1\leq i\leq n-1$ and there are no further comparabilities.
For $S=\{s\}$ and $C=P_n$, one has
$\mathcal W(S,C)=\{P_n\}$, so $(S,C)\in\Sat(P_n)$, and
\[
\bar\omega(S,C)
=
|\Max(P_n\setminus\{s\})|
=
n-1.
\]
Hence $\Omega(P_n)\geq n-1$.  Together with
\eqref{eq:uniform-global-dimension-bound} and
\Cref{thm:global-dimension-formula}, this gives
$\gldim\Lambda_{P_n}=n-1$.
\end{proof}


Finally, combining \Cref{thm:global-dimension-formula} with  \Cref{prop:relative-auslander} gives the
corresponding formula for the interval-resolution global dimension.

\begin{corollary}
\label{cor:interval-resolution-global-dimension-formula}
Let $P$ be a finite connected poset with $|P|>1$. Then
\begin{equation}
\label{eq:interval-resolution-global-dimension-formula}
\intresgldim_\kk P=\Omega(P)-2.
\end{equation}
If $|P|\geq3$, then
\begin{equation}
\label{eq:uniform-interval-resolution-bound}
\intresgldim_\kk P\leq |P|-3.
\end{equation}
\end{corollary}

\begin{proof}
The first equality follows from
\Cref{thm:global-dimension-formula} and \Cref{prop:relative-auslander}. The second follows from
\Cref{cor:uniform-global-dimension-bound}.
\end{proof}



